\section*{Executive summary}
\addcontentsline{toc}{section}{Executive summary}
\label{sec:execsummary}

We apply \emph{unbinned machine-learning unfolding}
(OmniFold~\cite{Andreassen:2019cjw, Andreassen:2021zzk}) to MINERvA's
medium-energy neutrino data. The analysis begins with a reproduction of a
published measurement and extends it to simultaneous three-, four-, and
five-observable scalar cross sections. We also study a recoil-only point-cloud
representation and a preliminary extended fiducial domain with relaxed
truth-muon cuts. The specific scalar extensions studied here do not have direct
binned MINERvA counterparts.

Traditional unfolding corrects detector effects with a migration matrix between
histogram bins, so each added observable multiplies the number of bins. OmniFold
instead learns an event weight that morphs simulation into data and transfers it
to the pre-detector level through the simulation's truth--reconstruction pairing.

Reporting binning is imposed only at the end. A recorded observable can therefore
be added without constructing a new response tensor, although training, support,
and uncertainty propagation still become harder with dimension.

The fixed-dimension measurements use gradient-boosted decision trees, validated
with closure, iteration, pull, and calibration diagnostics and neural-network
cross-checks. The existing 2D toy ensemble does not provide an independent-truth
coverage test; that test remains open. In five dimensions an end-to-end
independent-truth test was run on the statistical intervals of a deterministic
variant of the estimator, with a nominal-truth bias correction (1,200
pseudo-experiments over three truths). It failed its predeclared futility
criterion at the nominal and the $\Eavail$-tilted truths, so no frequentist
coverage is claimed for any five-dimensional interval.
% Source: docs/orchestration/OUTCOME-20260925-s5c-tier-s-futility-fail.md; VALIDATION_LEDGER VL150.
A successor development test replaced the per-bin purity weights by the unbinned
negative-weight background treatment. It removed most of the purity method's
nominal-truth bias in the highest-$W$ cells (from about \SI{-4}{\percent} to below
\SI{0.3}{\percent}), but left a bias of up to about \SI{1}{\percent} in some joint cells.
It also showed that the estimator itself does not follow a generator-sized change of the
$\Eavail$ shape: signal-only pseudo-experiments whose truth carries the GiBUU-to-GENIE
$\Eavail$ shape ratio come back several percent off in the largest $(\Eavail,W)$ cells
and up to \SI{74}{\percent} off in the highest-$W$ column. This regularization bias is
not in the quoted covariance.
% Source: docs/orchestration/OUTCOME-20260925-s5n-stage1-development-fail.md; VL151-VL152;
% KNOWN_ISSUES 75, 77.

\paragraph{Results overview:}
\begin{itemize}[itemsep=2pt]
  \item \textbf{Reproduction.} 
  Our unbinned measurement of the muon transverse and longitudinal momentum,  $d^{2}\sigma/(d\pT\,d\pPar)$, matches the published measurement~\cite{MINERvA:2021owq}: total cross-section ratio $\ratioTot$, \SI{\binsTen}{\percent} of bins within \SI{10}{\percent}, per-bin standardized differences sub-$1\sigma$, and a separately reconstructed uncertainty budget whose median per-bin uncertainty (\SI{\uqMedian}{\percent}) matches the published \SI{\uqPaper}{\percent} (\S\ref{sec:results}, \S\ref{sec:systematics}, Fig.~\ref{fig:xsec}).
  
  \item \textbf{Unbinned 3D/4D/5D MINERvA unfolds.}
  Adding available energy ($\Eavail$), momentum transfer ($q_{3}$), and hadronic mass ($W$) produces \emph{unbinned} simultaneous three-, four- and five-observable unfolds of MINERvA data; every result passes the projection anchors against its lower-dimensional predecessor to $\approx$1--2\,\% --- a built-in cross-check (\S\ref{sec:3d}). Prior MINERvA multi-differential cross sections, including a binned triple-differential quasielastic-like measurement~\cite{MINERvA:2022qe}, used explicit binned unfolding; the methodological distinction here is the unbinned, simultaneously-unfolded formulation rather than multi-dimensionality itself.
  
  \item \textbf{Established low-recoil discrepancy recovered.}
  All four truth-level generators (GENIE~\cite{Andreopoulos:2009rq}, MINERvA Tune~\cite{MINERvA:2016lowrecoil}, NuWro~\cite{Golan:2012wx}, GiBUU~\cite{Buss:2011mx}) under-predict the data at low $\Eavail$; pion final-state-interaction dials move the spectrum by less than \SI{1}{\percent} while adding Valencia 2p2h~\cite{Nieves:2011pp,Gran:2013kda} fills \SI{46}{\percent} of the data$-$CV gap in the quasielastic--$\Delta$ dip (\SI{27}{\percent} of the integrated deficit) --- the established MINERvA low-recoil/2p2h deficit~\cite{MINERvA:2016lowrecoil,MINERvA:2022incl}, re-found by an independent method (\S\ref{sec:3d-2p2h}, Fig.~\ref{fig:3dmodels}).
  
  The FSI statement concerns the tested pion dials only; it does not exclude
  other transport effects or uniquely identify the nuclear mechanism.

  \item \textbf{A new localization of the remaining excess.} 
  The high-$\Eavail$ excess that 2p2h cannot explain is localized in the $(\Eavail, W)$ plane: its positive data-minus-generator difference concentrates in the deep-inelastic corner ($W\ge\SI{1.8}{GeV}$), and all four generators underpredict that region. \textbf{A covariance for this projection \emph{is} adopted} --- the \zprojCells-cell projection \zprojDigest\ldots\ of the five-axis trunk \zcvDigest\ldots, adopted as publication-under-exception on 2026-09-20 and paired to the figure's central values by digest identity. The comparison is nonetheless reported at central-value level, and the reason is measured rather than precautionary: the estimator-seed sensitivity of that covariance is \SI{\sprojMeasured}{\percent} against a \SI{\sprojBound}{\percent} limit fixed before production, so no compatibility metric computed from it would state a precision the input has. That measurement and the three others that travel with the adoption are in \S\ref{sec:eavailw-adopted}; the adopted construction itself is specified in \S\ref{sec:eavailw-construction} (\S\ref{sec:eavailw}, Fig.~\ref{fig:eavailWband}).
  
  \item \textbf{A methodological finding.} 
  The standard covariance adds matrices built for one systematic source at a time, which omits cross-source response terms. In MINERvA's \texttt{GetTotalErrorMatrix}, each named source contributes its own covariance, built with all other sources held nominal; the sum retains within-source bin--bin correlations but contains no cross-source terms. The corrected unified-throw test perturbs all reweightable sources simultaneously using actual asymmetric endpoints, one fixed estimator seed, and mean-centered throws; it reports the joint mean shift separately. The selection-complete lateral replacement has since landed and one resulting 5D product is adopted as publication-under-exception, carrying four measurements that qualify everything built from it (\S\ref{sec:uq-combine}, \S\ref{sec:eavailw-adopted}); the other corrected candidates are not adopted.
  
  \item \textbf{Recoil-point-cloud representation cross-check (diagnostic, method development).}
  A point-cloud transformer uses reconstructed recoil clusters at step~1 and truth
  final-state hadrons at step~2; muon kinematics enter selection and binning but
  neither classifier. Its central value and ordinary closure therefore test a
  recoil-only estimator. The historical uncertainty assembly is quarantined
  because nuisance and retraining shifts were not combined with their cross term
  and detector membership was limited to CV support. A full-event estimator ---
  one trained jointly on the muon and recoil --- requires a new nominal fit,
  omitted-muon stress closure, and uncertainty budget (\S\ref{sec:pet},
  \S\ref{sec:pet-projection}). No PET
  covariance is adopted as a publication uncertainty product. The full-event
  central/statistical pairing was declined rather than resolved, with coverage
  the missing evidence.

  \item \textbf{A preliminary extended-fiducial study.}
  Re-running the scalar 2D $(\pT,\pPar)$ unfold with the truth-muon kinematic cuts removed enlarges the signal definition while retaining the tracker fiducial interaction volume and the unchanged MINOS-based reconstructed selection. About \SI{6}{\percent} of the rate lies in a newly acceptance-supported extension (tier~1), while about \SI{28}{\percent} lies in low- or zero-efficiency cells and is reported separately as prior-dominated extrapolation (tier~2). The result includes a restricted-region anchor, closure tests and a three-prior envelope, but not a corrected extended-fiducial covariance (\S\ref{sec:fps}).
 
  
\end{itemize}
